\documentclass[11pt]{article}
\usepackage{amsmath,amssymb,amsthm}
\newtheorem{lemma}{Lemma}[section]
\newtheorem{corollary}[lemma]{Corollary}
\newtheorem{definition}[lemma]{Definition}
\newtheorem{remark}[lemma]{Remark}
\newcommand{\WR}{W_R}
\newcommand{\WI}{W_I}
\newcommand{\ch}{\operatorname{ch}}
\newcommand{\Pic}{\operatorname{Pic}}
\newcommand{\KA}{A}
\newcommand{\KAd}{\widehat{A}}

\title{GST Picard--Sheaf Flatness of the Pair Channel: the Entry Wall of the
Monodromy Mechanism (Lemma W8) --- Round 15-b}
\author{GST Hodge campaign, cycle 4 follow-up (SKILL V6 MODE B, ULTRA)}
\date{2026-10-26}

\begin{document}
\maketitle

\section*{Scope box (read first)}

\textbf{Claimed here:} Lemma W8 (every holomorphic pure pair on the
discriminator has a FLAT determinant-of-cohomology difference line bundle on
the dual fourfold); the monodromy observable $\Psi_{\mathrm{pair}}$;
Corollaries W8.1 (the entry-test extension) and W8.2 (the FM-closure
consistency at the entry layer). \textbf{NOT claimed:} any decision of the
dichotomy bit; any computation of a monodromy representation; any rigidity;
any exhaustiveness; any sheaf-level statement on the dual side beyond the
committed WIT status; any resolution of the campaign contract. This strike
\textbf{DETERMINES NO BIT}; it opens the registry's named NO-route reopen
mechanism (Picard--sheaf monodromy / determinant-of-cohomology rigidity) with
a certified entry wall and its observable, and records the mechanism's live-or-dead
conditions. Labels W8/W8.1/W8.2 extend the committed W-frame; no committed
label is reused.

\section{Setting and inputs}

$\KA$ is the committed nonsplit $\mathbb{Q}(i)$-Weil principally polarized
fourfold discriminator (NS-rank one, $\theta$ the polarization class,
$\WR,\WI$ the Weil classes; committed constants
$\int\theta^4=24$, $\int \WR^2 = \int \WI^2 = 8$, $\int \WR\WI = 0$, all
$\int \WR\,\theta^j = 0$). $\KAd$ is the dual abelian fourfold.
A \emph{holomorphic pure pair at scale $m\ge 1$} is the committed
pair-form input: actual holomorphic vector bundles $E_\pm$ on $\KA$ with
equal ranks and
\[
\ch(E_+) - \ch(E_-) \;=\; m\,\WR
\]
\emph{purely at grade 2}: the grade-0, grade-1, grade-3 and grade-4
components of the difference all vanish. (By the committed W6 table the
members then have equal $c_1 \in \mathbb{Z}\theta$, equal $c_3$ --- the
$c_3$-shift $-m\,c_1\wedge \WR$ dies by the W0 annihilation
$\theta^j \wedge \WR = 0$ as a class --- a $c_2$-difference $-m\,\WR$,
and Weil charges differing by exactly $-8m$.) The existence of one such
pair at any scale is equivalent to the vertex bit $S = W_\Q$ (the committed
W2/W1 equivalence, M3-normalized). Both polarities remain armed; the
registry reopen menu names, for the NO route, ``a non-additive, FM-closed
rigidity mechanism (refined Brill--Noether-type data, Picard--sheaf
monodromy, determinant-of-cohomology rigidity)''. This strike opens that
mechanism's entry layer.

\paragraph{Committed inputs used.} The pair definition and the W6 table
(kernel-green, run 40139821467 @ c6a2778); the W0 class annihilation
(kernel-green, same run); the W5 cohomological FM normalization
(kernel-green, same run; the exact dual-basis convention, coefficient 1,
both round trips the identity); the W4 difference-Euler corridor
(kernel-green, same run).

\paragraph{External leaves (declared, standard strength; not certified
here).}
\begin{itemize}
\item[P1] \emph{Determinant--K-theory identity:} for $\xi \in K^0(X)$
represented by a perfect complex, $\det(\xi) \in \Pic(X)$ is defined and
satisfies $c_1(\det \xi) = \ch_1(\xi)$. (Standard: on vector bundles
$c_1(\det E) = c_1(E) = \ch_1(E)$; $\det$ is a group homomorphism
$K^0 \to \Pic$ and $\ch_1$ is additive.)
\item[P2] \emph{GRR for the Poincar\'e family on an abelian variety:}
with $\mathcal{P}$ the normalized Poincar\'e bundle on $\KA\times\KAd$ and
$p_{\KAd}$ the projection,
$\ch(Rp_{\KAd *}(p_\KA^*F\otimes m^*\mathcal{P}))\cdot \mathrm{td}(\KAd)
= p_{\KAd *}(\ch(F)\, e^{c_1(\mathcal{P})}\,\mathrm{td}(\KA)\,
\mathrm{td}(\KAd))$; and $\mathrm{td}(\KA) = \mathrm{td}(\KAd) = 1$
because the tangent bundle of an abelian variety is holomorphically
trivial (all its Chern classes vanish). Hence
$\ch(\widehat{\Phi}(F)) = p_{\KAd *}(\ch(F)\, e^{c_1(\mathcal{P})})$.
\item[P3] \emph{The cohomological FM grade reflection,} extended to all
grades by the same committed construction: with the exact dual-basis
convention $c_1(\mathcal{P}) = \sum_{p=0}^{7} u_p \otimes \widehat u_p$,
the component of $\ch(\widehat\Phi(F))$ in $H^{2(4-k)}(\KAd)$ is the
pairing of $\ch_k(F)$ against $c_1(\mathcal{P})^{\,8-2k}/(8-2k)!$
(the pushforward along the first factor). At $k=2$ this is exactly the
committed kernel-green W5 transform; at $k=3$ it gives
$\ch_1(\widehat\Phi(F)) = \mathrm{fm}_3(\ch_3(F))$, a degree-6 to
degree-2 word map of the same exterior model (census check: dimensions
$28 \to 28$, both sides $=\binom{8}{6}=\binom{8}{2}$).
\item[P4] \emph{Duality bookkeeping:} $\Pic^0(\KAd) \cong
\widehat{\widehat{A}} \cong \KA$ canonically (standard).
\end{itemize}

\section{Lemma W8 and its consequences}

\begin{lemma}[W8: Picard--sheaf flatness of the pair channel]
\label{lem:W8}
Let $(E_+,E_-)$ be a holomorphic pure pair on the discriminator $\KA$ at any
scale $m\ge 1$ (any rank $r\ge 1$). Write
\[
\lambda(F) \;:=\; \det\bigl(Rp_{\KAd *}(p_\KA^*F\otimes m^*\mathcal{P})\bigr)
\;\in\; \Pic(\KAd)
\]
for the determinant-of-cohomology line bundle, and
\[
L_{\mathrm{pair}} \;:=\; \lambda(E_+)\otimes \lambda(E_-)^{-1}\;\in\;\Pic(\KAd).
\]
Then
\[
c_1\bigl(L_{\mathrm{pair}}\bigr) \;=\; 0.
\]
Since $L_{\mathrm{pair}}$ is a line bundle, all higher Chern classes vanish,
and $c_1 = 0$ means $L_{\mathrm{pair}} \in \Pic^0(\KAd)$: the pair's
determinant-of-cohomology difference is a FLAT line bundle on the dual
fourfold.
\end{lemma}

\begin{proof}
Let $\xi_{\mathrm{pair}} := [E_+] - [E_-] \in K^0(\KA)$. The perfect-complex
determinant is a group homomorphism $\det : K^0(\KAd) \to \Pic(\KAd)$ (leaf
P1), and the FM transform is exact, so
$[\widehat\Phi(\xi_{\mathrm{pair}})] = [\widehat\Phi(E_+)] -
[\widehat\Phi(E_-)]$ in $K^0(\KAd)$ and therefore
\[
L_{\mathrm{pair}}
= \det\bigl(\widehat\Phi(\xi_{\mathrm{pair}})\bigr).
\]
Using leaf P1 for the first equality and leaf P3 (the grade bookkeeping of
leaf P2: the grade-1 component of the transform receives exactly the
grade-3 component of the input, $k=3 \mapsto 4-k=1$) for the second,
\[
c_1(L_{\mathrm{pair}})
= \ch_1\bigl(\widehat\Phi(\xi_{\mathrm{pair}})\bigr)
= \mathrm{fm}_3\bigl(\ch_3(\xi_{\mathrm{pair}})\bigr).
\]
By the pair definition, $\ch(\xi_{\mathrm{pair}}) = m\,\WR$ purely at
grade 2, so $\ch_3(\xi_{\mathrm{pair}}) = 0$. Linearity of $\mathrm{fm}_3$
gives $c_1(L_{\mathrm{pair}}) = \mathrm{fm}_3(0) = 0$.
\end{proof}

\begin{remark}[Where the walls sit inside the proof]
The proof consumes the pair definition's grade-3 purity and nothing else
from the wall system; the W6 equal-$c_3$ table and the W0 corridor are the
certified \emph{mechanism} showing that purity at grade 3 is compatible
with actual Chern data (the $c_3$-shift $-m\,c_1\wedge\WR$ dies by W0 in
the ring). No NS-rank input enters the proof of W8 itself: the flatness is
a formal consequence of pair purity $+$ abelianness $+$ the determinant
identity, and therefore holds word-for-word on every abelian world,
including the split worlds (the wall is world-uniform as a necessary
condition --- the standard consistency-probe pattern; it kills nothing by
itself).
\end{remark}

\begin{definition}[the monodromy observable]
\label{def:Psi}
For a holomorphic pure pair $(E_+,E_-)$ on $\KA$,
\[
\Psi_{\mathrm{pair}}(E_+,E_-) \;\in\; \Pic^0(\KAd) \cong \KA
\]
is the point of the double-dual (leaf P4) representing the flat line
$L_{\mathrm{pair}}$. $\Psi_{\mathrm{pair}}$ is defined for every
holomorphic pure pair; it is the entry observable of the Picard--sheaf
monodromy mechanism. $\Psi_{\mathrm{pair}}$ is NOT claimed to be
non-additive here; whether the observable (together with the monodromy
representations of the per-fiber virtual objects
$R\Gamma(E_\pm\otimes P)$, $P\in\Pic^0(\KA)$) constrains the charge data
$(m,q)$ beyond the additive profile is exactly the mechanism's open
question (see \S\ref{sec:live-or-dead}).
\end{definition}

\begin{corollary}[W8.1: the entry-test extension]
\label{cor:W81}
Any claimed holomorphic pure pair on the discriminator --- equivalently any
claimed W-charged surface output (the committed W2/W1 equivalence) --- must
now clear one further entry test beyond the committed list
($\chi$-corridor blindness, no cohomological band, the W6 table, the End-mass
bound, FM-closure): it must exhibit its flat line
$L_{\mathrm{pair}}\in\Pic^0(\KAd)$ and the monodromy representation of the
per-fiber virtual FM objects. The flatness $c_1 = 0$ is a new necessary
condition, certified at the model layer for the grade bookkeeping and at
leaf strength for the scheme layers P1/P2.
\end{corollary}

\begin{corollary}[W8.2: the FM-closure consistency at the entry layer]
\label{cor:W82}
The entry layer passes the W5 duality requirement at the K-theoretic level:
the dual-side difference $\widehat\Phi(\xi_{\mathrm{pair}})$ has
$\ch = m\,\mathrm{fm}_2(\WR) = m\,\widehat\WR$-dual (the committed W5
normalization, kernel-green, coefficient exactly 1) at grade 2 and
vanishing grades $0,1,3$ --- the same purity shape on $\KAd$ --- so the
double-dual layer's flatness holds identically, and the mechanism is
duality-symmetric at the level where it is defined. No WIT-per-point claim
is made: the per-fiber FM of a pair is a complex with cancelling Euler by
the committed W4 corridor (the difference-Euler vanishes against every
$P\in\Pic^0$), and that remains its certified status.
\end{corollary}

\section{Live-or-dead conditions for the mechanism}
\label{sec:live-or-dead}

The registry reopen rule requires each armed mechanism to carry its own
closure conditions. For the Picard--sheaf monodromy mechanism:

\begin{itemize}
\item \textbf{LIVE} (a non-additive FM-closed wall exists): if the monodromy
representations of the per-fiber virtual objects
$R\Gamma(E_\pm\otimes P)$ as $P$ winds around $\Pic^0(\KA)$, together with
the flat position $\Psi_{\mathrm{pair}}$, impose a constraint on the charge
data $(m,q)$ that the additive profile does not compute --- and that
constraint is incompatible with one of the two strata --- then the mechanism
separates the strata and the NO route reopens with a certified wall.
\item \textbf{DEAD} (the mechanism closes): if both strata are realized by
monodromy-consistent data (the analog of the additive-completeness
consistency), the mechanism dies and the registry row records the closure,
exactly as the six NO-route attack lines of cycle 4 died.
\end{itemize}

Neither outcome is claimed here. The certified content of this round is the
entry wall (W8), the observable, the entry-test extension, and the
duality consistency --- the mechanism's first certified step from bare menu
item to armed route.

\section{Verification layer}

\texttt{research/experiments/task15a\_picard\_flatness\_checks.py}
(committed alongside this dossier). Bounds: pure Python, exact
\texttt{Fraction} arithmetic throughout, no floats, zero tolerance, every
check ASSERTS; the committed 16-letter exterior model conventions imported
verbatim from the audited task13d certificate (parity-exact Weil
construction, $\theta$, the exact dual-basis Poincar\'e class, the Newton
conductor machinery, the committed instance data); largest enumeration
$\le 2^{16}$ monomial operations; runtime $< 1$\,s. It verifies:
(1) the committed constants re-asserted (periods, W0 annihilations ---
the inherited-artifact gate re-run per ledger 084);
(2) the general-grade FM bookkeeping: the $\mathrm{fm}_k$ map
($k=1,2,3$) by the identical committed construction
$c_1(\mathcal{P})^{8-2k}/(8-2k)!$, with the $k=2$ case verified IDENTICAL
to the committed W5 transform on $\WR,\WI$ (coefficient 1, round trips);
(3) the $k=3$ census: degree-6 words map to degree-2 words
($28 \to 28$), non-vacuity asserted ($\mathrm{fm}_3(\theta^3)\ne 0$),
and the composite $\widehat{\mathrm{fm}}_1\circ \mathrm{fm}_3$ asserted
uniform on a spanning test set with constant of modulus 1 (the
normalization constant of the committed convention at grade 3, recorded
by the script, not fitted);
(4) the pair grade-3 purity re-derived on ALL committed instances (the
four W6 instances and the carrier instance) through the conductor
machinery: $\ch_3(E_+) - \ch_3(E_-) = 0$ and the $c_3$-shift
$-m\,c_1\wedge\WR$ dying by W0 in the ring when $c_1\in\mathbb{Z}\theta$;
(5) the flatness composites $\mathrm{fm}_3(\ch_3\text{-difference}) = 0$
on all instances (the W8 identity at the model layer);
(6) the duality face: $\mathrm{fm}_2(m\,\WR) = m\,\widehat\WR$-dual exact
(linearity, coefficient $m$, W5-consistent) for W8.2.

The scheme-level steps (leaves P1, P2, P4 and the perfect-complex
determinant theory) are declared external leaves at standard strength and
are NOT claimed as certified; the model layer certified here is the same
layer as the committed W5/W6 certificates (kernel-green run 40139821467 @
c6a2778), and no new Lean file is introduced by this round (the kernel
content of W8's model layer --- pair purity at grade 3 --- is already
asserted inside the committed walls certificate; this is receipt-cited, not
re-certified).

\section{Provenance}

Builds on: the cycle-4 strike-1 walls (W0/W4/W5/W6, kernel-green run
40139821467 @ c6a2778, audit-closed 14-b-3) and the cycle-4 strike-2
decision attempt (W7, the funnel, the additive-completeness theorem;
audit-closed 14-b-3); the round 15-a terminal-state exactness repair
(ledger 085, commit 26214ce). The registry NO-route reopen item
``Picard--sheaf monodromy / determinant-of-cohomology rigidity'' is the
executed route. Research-only round: no .lean, no .yml, no workflow fires.

\end{document}
